% !TeX root = ../main.tex
\section{Reinforcement Learning in Imperfect-Information Games}\label{sec:rl}

The final rung of the algorithmic ladder replaces equilibrium computation with \emph{trial and error}: an agent that learns by playing, observing rewards, and updating its own policy. Reinforcement learning (RL) is the most flexible of the paradigms surveyed here---it needs no game-tree enumeration, no explicit opponent model, and no hand-crafted thresholds---and consequently faces the sharpest theoretical hazards. This section sets out the classical RL framework, explains precisely where poker breaks its assumptions, describes the self-play designs that restored soundness (culminating in DeepStack, Libratus, and Pluribus), and specifies our own tabular Q-learning agent, which shares the CFR agents' abstraction so that Section~\ref{sec:experiments} can attribute performance differences to the learning paradigm rather than the state representation.

\subsection{Markov decision processes and Q-learning}

The classical substrate is the Markov decision process (MDP): states $s \in \mathcal{S}$, actions $a$, transition kernel $P(s' \mid s, a)$, rewards $r(s,a)$, and discount $\gamma \in [0,1)$. The Bellman optimality equation characterises the optimal action-value function:
\begin{equation}\label{eq:bellman}
Q^{*}(s,a) \;=\; r(s,a) + \gamma\, \E_{s' \sim P(\cdot \mid s,a)}\!\left[\max_{a'} Q^{*}(s',a')\right].
\end{equation}
Tabular Q-learning iterates the stochastic approximation
\begin{equation}\label{eq:qlearn}
Q(s_t, a_t) \;\leftarrow\; (1 - \alpha_t)\, Q(s_t, a_t) \;+\; \alpha_t \left[ r_t + \gamma \max_{a'} Q(s_{t+1}, a') \right],
\end{equation}
whose target $r_t + \gamma \max_{a'} Q(s_{t+1},a')$ is itself a biased estimate built from the current $Q$---a \emph{bootstrapped} update. Watkins' theorem guarantees convergence $Q \to Q^{*}$ under standard stochastic-approximation conditions: every state--action pair visited infinitely often (sufficient exploration), and step sizes $\alpha_t$ decaying appropriately~\cite{watkins1992q}. In a poker engine the episode structure is natural: one hand of play is an episode, the terminal reward is the chip delta, and \eqref{eq:qlearn} is propagated backwards over the seat's decision sequence within the hand.

\subsection{Where the substrate breaks}

Poker violates every MDP assumption in turn. The state is not observed: the transition ``kernel'' depends on opponents' hidden cards, so the learner's information state is a belief, not a state, and the process over information states is not Markov in general. The environment is not stationary: opponents adapt, so the very dynamics being estimated drift as the learner improves---the condition $\alpha_t \to 0$ with infinite visitation becomes meaningless when the thing visited keeps changing. And the objective itself bifurcates: maximising \eqref{eq:bellman} against a \emph{fixed} environment is exploitation, whereas Section~\ref{sec:gametheory} established that adversarial settings reward unexploitability. An RL agent trained against one opponent pool learns to beat that pool and may thereby learn habits that a different opponent punishes; the analogous failure of the rule-based agent (Section~\ref{sec:evaluation}) was static thresholds, and the cure---equilibrium reasoning---is the same.

The classical remedy is \emph{self-play}: let the environment be other learners of the same family. Fictitious play, the grandfather of self-play schemes, has each player best-respond to the time-average of its opponents' past play; Robinson proved convergence to equilibrium in two-player zero-sum games~\cite{robinson1951iterative}. Self-play with function approximation underlies the landmark game-playing results---TD-Gammon's surprise strength in backgammon~\cite{tesauro1995td} and the AlphaGo/AlphaZero line in perfect-information games~\cite{silver2017alphagozero}. In imperfect-information games, however, naive deep self-play can converge to brittle, mutually overfit policies, which motivated the hybrid architectures below.

\paragraph{NFSP.} Neural fictitious play (NFSP) makes Robinson's construction scalable: each agent carries \emph{two} policies, a best response learned by deep Q-learning against the recent opponent distribution, and an average policy learned by supervised regression over the agent's own history of best-response play~\cite{heinrich2016deep}. Playing the average policy approximates fictitious play (hence equilibrium convergence in two-player zero-sum), while periodically best-responding supplies the exploitation pressure that keeps the average honest. NFSP was the first RL-based agent to approach competitive strength in heads-up limit Hold'em without explicit equilibrium computation, and its two-headed design---exploit \emph{and} average---remains the template for reconciling RL with game theory.

\paragraph{The deadly triad.} Wherever function approximation meets bootstrapping and off-policy data, value iteration can diverge---Sutton and Barto's ``deadly triad''~\cite{sutton2018rl}. Deep self-play adds a fourth hazard: the data distribution shifts as all policies move. Production mitigations (target networks, replay reservoirs, trust regions) are engineering responses to instability rather than fixes for it; our miniature deliberately stays tabular, accepting a small state space instead of an unstable one, and quantifies the cost of that choice in Section~\ref{sec:experiments}.

\subsection{The superhuman architectures}

\paragraph{DeepStack (2017).} DeepStack's designers rejected the blueprint-plus-search template and instead \emph{continually re-solved} the remaining subgame at every decision~\cite{moravcik2017deepstack}. A re-solve starts from the current belief over the opponent's holdings (continual re-solving: the belief is inherited from the previous re-solve's strategy, not recomputed), solves the depth-limited continuation of the game tree by CFR-style iteration, and evaluates \emph{leaf} information sets with a deep network trained offline on millions of exactly solved subgames. The soundness argument is the paper's theoretical contribution: because each leaf evaluation enters the solved value linearly, the strategy's exploitability is bounded by the leaf-network's error weighted by its influence---unlike heuristic forward search, the procedure degrades gracefully and provably. DeepStack's second network evaluated full-hand distributions preflop, and its cascades of re-solves defeated a panel of human professionals in heads-up no-limit Hold'em.

\paragraph{Libratus (2017).} Libratus is the archetype of the blueprint-plus-search school~\cite{brown2017libratus}. Its blueprint---an equilibrium of a coarse abstraction computed by external-sampling MCCFR with CFR$^{+}$-style updates over 64-core months of self-play---played the early streets. From the third betting round onward, Libratus \emph{re-solved} the actual subgame in real time on the unabstracted betting sequence, using a nested safety condition to guarantee the re-solved strategy is no more exploitable than the blueprint's recommendation. Between sessions, a \emph{self-improver} tightened the endgame abstraction where its human opponents had concentrated their play (identified via their shown-down hole cards), so the system strengthened overnight against the specific humans it was fighting. Libratus won its 2017 bridge match against four human specialists by 1.76 million chips over 120{,}000 hands---a margin far beyond statistical doubt.

\paragraph{Pluribus (2019).} Pluribus carried the template to six-player no-limit Hold'em, where equilibrium guarantees dissolve (Section~\ref{sec:nash}) and compute budgets matter~\cite{brown2019pluribus}. Its blueprint cost about $\$150$ of cloud compute---three orders of magnitude less than AlphaGo's hardware---because it used a cheap ``bubble'' abstraction and a modest MCCFR budget. Live play used \emph{limited-lookahead search}: within the current betting round, Pluribus enumerated unabstracted betting sequences a few actions deep and evaluated leaf information sets with the blueprint's continuation values, with a ``sticky'' variant of card abstraction to keep leaf evaluations consistent across the search. Against five humans at once, Pluribus won at rates professional players described as unsurprising for a top human specialist; the AI rating systems of the day considered it superhuman. The lesson Pluribus adds to Libratus is architectural parsimony: equilibrium-quality search, applied surgically at the decision points that matter, converts a modest blueprint into a superhuman whole.

\subsection{Our RL agent: tabular Q-learning with an opponent pool}

Our implementation instantiates the paradigm in its honest, small-scale form. The agent learns over \emph{exactly the same information-set key space} as the CFR agents---position class, street, card bucket, betting pattern---so that the experimental comparison of Section~\ref{sec:experiments} isolates the learning rule. Training is tabular Q-learning \eqref{eq:qlearn} in a six-seat self-play loop in which the learner occupies one seat and a \emph{mixed opponent pool} occupies the other five: three rule agents, one uniform-random agent, and one CFR-blueprint agent. The pool's composition is a deliberately cheap stand-in for NFSP's averaging: facing a stationary mixture of heterogeneous opponents keeps the target of \eqref{eq:qlearn} approximately stationary while exposing the learner to value-oriented, erratic, and equilibrium-flavoured opposition in fixed proportion. Rewards are the terminal chip delta of each hand, discounted with $\gamma = 0.9$ and propagated backwards over the learner's decision sequence; exploration is $\varepsilon$-greedy with $\varepsilon$ annealed during training (and held at $0.02$ in play). No poker knowledge is hand-coded: aggression thresholds, positional discipline, and pot-odds awareness emerge, or fail to emerge, from the reward signal alone.

Three limitations should be stated plainly, because Section~\ref{sec:experiments} measures their consequences. First, Q-learning here optimises \emph{exploitation of the pool}, not equilibrium play: against the specific five-opponent mixture the learned policy is near-optimal, but nothing protects it outside that mixture---the guarantee asymmetry of Section~\ref{sec:nash} in miniature. Second, the tabular representation ties the agent to the shared abstraction, so it inherits the abstraction's blind spots without Deep CFR's ability to generalise across them. Third, the discounted-terminal-reward formulation teaches the agent hand-level chip maximisation only: it cannot represent multi-hand strategic effects (image, trapping across hands), which the superhuman systems likewise ignored and which poker folklore overrates. Within these limits the agent provides the experiment's cleanest test of a single question: \emph{how far does pure trial and error go, in a game where the environment pushes back?}
